\section{Multivariate beta and fixed-dimensional simplex integrals}
\label{sec:multivariate}

For a finite nonempty list
$\boldsymbol\alpha=(\alpha_1,\ldots,\alpha_n)$ with positive real parts, define
\[
 \B(\boldsymbol\alpha)
 =\frac{\prod_{j=1}^{n}\Gamma(\alpha_j)}
        {\Gamma(S)},
 \qquad S=\sum_{j=1}^{n}\alpha_j.
\]
For $n=2$ this is the ordinary beta function.

\subsection{Coordinate recurrences and Pascal's identity}

\begin{theorem}[Multivariate recurrence]\label{thm:multi-recurrence}
For each $j$,
\[
 \boxed{
 \B(\alpha_1,\ldots,\alpha_j+1,\ldots,\alpha_n)
 =\frac{\alpha_j}{S}\B(\boldsymbol\alpha).}
\]
Consequently
\[
 \boxed{
 \B(\boldsymbol\alpha)=
 \sum_{j=1}^{n}
 \B(\alpha_1,\ldots,\alpha_j+1,\ldots,\alpha_n).}
\]
\end{theorem}

\begin{proof}
Apply $\Gamma(z+1)=z\Gamma(z)$ in the numerator and denominator.  Summing the coordinate formulas multiplies $\B(\boldsymbol\alpha)$ by
$\sum_j\alpha_j/S=1$.  For internalization, clear the common denominator $S$ before the finite sum.  The list sum and product are primitive-recursive recursors.
\end{proof}

Repeated shifts give, for a vector $\mathbf m\in\N^n$,
\[
 \B(\boldsymbol\alpha+\mathbf m)
 =
 \frac{\prod_j(\alpha_j)_{m_j}}
      {(S)_{\sum_jm_j}}
 \B(\boldsymbol\alpha).
\]
At positive integer arguments this becomes
\[
 \B(m_1,\ldots,m_n)
 =\frac{\prod_j(m_j-1)!}
        {(\sum_jm_j-1)!}.
\]
This is the direct analogue of the multinomial coefficient relation.

\subsection{Derivative formula}

The source-page derivative
\[
 \frac{\partial}{\partial\alpha_j}\B(\boldsymbol\alpha)
 =
 \B(\boldsymbol\alpha)\bigl(\psi(\alpha_j)-\psi(S)\bigr)
\]
was proved in \eqref{eq:multibeta-derivative}.  Higher derivatives are again finite Bell polynomials in polygamma differences.  For each fixed order and list length, this is a finite formal identity.  A variable list can be represented by a code and a length; the internal theorem quantifies only over the code and uses bounded lookup.

\subsection{Simplex integral}

For fixed $n\ge2$, let
\[
 \Delta_{n-1}=
 \{(t_1,\ldots,t_{n-1})\in[0,1]^{n-1}:
   t_1+\cdots+t_{n-1}\le1\},
 \qquad
 t_n=1-\sum_{j<n}t_j.
\]
Then
\begin{equation}
 \boxed{
 \int_{\Delta_{n-1}}\prod_{j=1}^{n}t_j^{\alpha_j-1}
 \,dt_1\cdots dt_{n-1}
 =\B(\boldsymbol\alpha).}
 \label{eq:simplex-beta}
\end{equation}
A constructive proof avoids a general multivariate integration theorem.  Integrate $t_{n-1}$ first and use the substitution
\[
 t_{n-1}=(1-t_1-\cdots-t_{n-2})u.
\]
The inner integral becomes
\[
 (1-t_1-\cdots-t_{n-2})^{\alpha_{n-1}+\alpha_n-1}
 \B(\alpha_{n-1},\alpha_n).
\]
Repeat.  After $n-2$ steps, the product of beta factors telescopes to the Gamma quotient.  Every stage is a one-dimensional finite Riemann sum with endpoint bounds; the dimension $n$ is externally fixed, or encoded by a simultaneous primitive recursion if a uniform family is desired.

\subsection{Dirichlet normalization without probability axioms}

For positive real parameters, dividing the simplex integrand by
$\B(\boldsymbol\alpha)$ produces a finite-sum name whose integral is $1$.  Expectations of monomials are merely ratios of shifted multivariate beta values:
\[
 \frac{1}{\B(\boldsymbol\alpha)}
 \int_{\Delta_{n-1}}
 \prod_j t_j^{\alpha_j+m_j-1}\,dt
 =
 \frac{\prod_j(\alpha_j)_{m_j}}
      {(S)_{\sum_jm_j}}.
\]
Thus the algebraic moment formulas usually stated for the Dirichlet distribution are already iPRA theorems about finite weighted sums and their certified limits; no measure-theoretic probability space is required.
